% !PRINTED 708
\weberpage{708}
Man erhält die Zahlen $n$, wenn man von der Gesammtheit aller natürlichen
Zahlen $k$ die Vielfachen von $q$, d.\,h. die Gesammtheit der Zahlen $qk$
wegnimmt. Hiernach ist
\[
\Sigma \frac{1}{n^s} = \Sigma \frac{1}{k^s} - \Sigma \frac{1}{q^s k^s} = (1-q^{-s}) \Sigma \frac{1}{k^s}.
\]
Wenn man aber in dem Satze 3., \S.\,191, für die $\mu_n$ die Reihe der
natürlichen Zahlen $k$ setzt, so folgt:
\[
\underset{s=1}{\mathrm{Lim}} (s-1) \Sigma \frac{1}{k^s} = 1,
\]
und wir erhalten daher
\weq{12}{\underset{s=1}{\mathrm{Lim}} \frac{s-1}{1-q^{-s}} \Sigma \frac{1}{n^s} = 1.}
Die anderen Factoren des Productes $\Phi(s)$ aber sind, wenn die
unendlichen Reihen nach steigenden Werthen von $n$ geordnet werden, nach
dem Satze 1., \S.\,191, stetige Functionen von $s$.

Denn die nach steigenden Werthen von $n$ geordnete Summe
\weq{13}{\Sigma \chi(n)}
ist zwar nicht convergent, kann aber doch dem absoluten Werthe nach nicht
über eine endliche Grenze hinausgehen. Denn so oft $n$ ein volles
Restsystem nach dem Modul $m$ durchläuft, kommt zu der Summe (13) ein
Beitrag hinzu, der nach \S.\,11, 6. den Werth 0 hat. Demnach ist, wenn $\chi$
nicht der Hauptcharakter ist,
\weq{14}{\underset{s=1}{\mathrm{Lim}} \overset{n}{\Sigma} \frac{\chi(n)}{n^s} = \overset{n}{\Sigma} \frac{\chi(n)}{n},}
und demnach erhalten wir für die Classenzahl $h$ im Kreistheilungskörper
$\Omega_m$ jetzt den Ausdruck
\weq{15}{g\,h = \overset{\chi}{\Pi} \overset{n}{\Sigma} \frac{\chi(n)}{n},}
in dem sich das Productzeichen $\Pi$ nur noch auf die \gesperrt{vom Hauptcharakter
verschiedenen} Charaktere $\chi$ bezieht.

\begin{center}
\S.\,194.

\gesperrt{Bestimmung der Summen X.}
\end{center}

Die unendlichen Reihen
\weq{1}{X = \overset{n}{\Sigma} \frac{\chi(n)}{n},}
%CONT>>
